\documentclass[exercice,a4]{thedocument}%
\usepackage{pstricks}%
\usepackage{tabularx}%
\usepackage[np]{numprint}%
\usepackage{enumitem}%
\startdatakeys
    \source{25GENMATAN1}
    \BO{2026}
    \cycle{4}
    \competencesDuSocle{CA,RA}
    \objectifsApprentissage{C4A1aa,C4A2ad,C4B2cb,C4D2aa,C4D3ac}
    \motsCles{probabilité,facteurs premiers,pourcentage,agrandissement,moyenne}
\enddatakeys
\begin{document}%
    
    \def\euro{\texteuro}%

\emph{Dans cet exercice, les cinq situations sont indépendantes. Il est rappelé que chaque réponse doit être justifiée sauf indication contraire.}

\medskip
\begin{itemize}[label=$\bullet~~$]
\item \textbf{Situation 1}

Dans une urne de $40$ boules indiscernables au toucher, $5$ sont rouges, $20$ sont vertes et $15$ sont blanches. L'expérience consiste à tirer au hasard une boule de l'urne et à noter sa couleur.

\smallskip

Calculer la probabilité d'obtenir une boule verte.
\item \textbf{Situation 2}

Décomposer en produit de facteurs premiers le nombre \np{1050}.

\emph{Aucune justification n'est attendue}.
\item \textbf{Situation 3}

Un article coûte $25$~\euro. Calculer son prix après une augmentation de $14\,\%$.
\item \textbf{Situation 4}

\begin{tabularx}{\linewidth}{X|X}
Le polygone 2 est un agrandissement du polygone 1.

Le coefficient de cet agrandissement est $2,5$.

L'aire du polygone 1 est égale à $7,5$~cm$^2$.

Calculer l'aire du polygone 2.&\emph{La figure ci-dessous n'est pas à l'échelle}.

\begin{center}
\psset{unit=0.7cm}
\begin{pspicture}(0,-1)(6,4)
\pspolygon(0,0)(3,0)(3,2)(1,3)(0,2)
\pspolygon(4.7,0.7)(5.7,0.7)(5.7,1.4)(5.1,1.7)(4.7,1.4)
\rput(1.5,-1){Polygone 2}\rput(5.2,-1){Polygone 1}
\end{pspicture}
\end{center}
\end{tabularx}
\item \textbf{Situation 5}

\medskip

Dans une classe de 3\up{e} on note la répartition des tailles des élèves dans le tableau suivant:

\begin{tabularx}{\linewidth}{|m{1.25cm}|*{8}{>{\centering \arraybackslash}X|}}\hline
\textbf{Taille (en cm)}&152&157& 160& 162& 165&170& 174& 180\\ \hline
\textbf{ Effectif}& 2& 4& 2& 5& 2&4&6&5\\ \hline
\end{tabularx}

\medskip

\begin{enumerate}
\item Quelle est la moyenne des tailles des élèves de cette classe ?
\item Quelle est la médiane des tailles des élèves de cette classe ?
\end{enumerate}
\end{itemize}

\newpage
\showthesource
\end{document}
